\documentclass[11pt]{article}
\usepackage[margin=.76in]{geometry}
\usepackage{amsmath,amssymb}
\usepackage[colorlinks=true,linkcolor=black,urlcolor=blue]{hyperref}
\newcommand{\tr}{\operatorname{Tr}}
\newcommand{\av}{\mathcal A_{n,m}}
\title{\vspace{-2em}The refined BMV trace inequalities\\for a matrix with two spectral values}
\hypersetup{pdfauthor={GPT-6.2 Astra (high), Aryan Singh},pdftitle={The refined BMV trace inequalities for a matrix with two spectral values}}
\author{GPT-6.2 Astra (high)\enspace{\small(lead AI author of this edition)}\\\href{https://www.researchgate.net/profile/Aryan-Singh-101}{Aryan Singh}\enspace{\small$\cdot$\enspace oo research\enspace$\cdot$\enspace The Open University (student)}}
\date{26 November 2026}
\begin{document}
\maketitle
\begin{abstract}
Let $A,B$ be positive semidefinite matrices of any finite size. If $A$ has
at most two distinct eigenvalues, the mean trace of all words with $n$
letters $A$ and $m$ letters $B$ lies between $\tr[(A^{n/m}B)^m]$ and
$\tr(A^nB^m)$ for all positive integers $n,m$. This follows by grouping
H\"agele's path expansion according to the number of visits to one eigenspace.
The resulting weights are coefficients of an ordinary Bessis--Moussa--Villani
polynomial and are therefore nonnegative. Together with a known
three-eigenvalue counterexample, this gives the sharp spectral threshold
for the upper bound. An appendix independently reproduces finite cases
of a previously established three-eigenvalue lower bound.
\end{abstract}

\section*{The two-sided bound}
Write $\av(A,B)=\binom{n+m}{m}^{-1}\sum_W\tr W$, where $W$ runs over all
words with $n$ letters $A$ and $m$ letters $B$. H\"agele's refinement of the
BMV conjecture asks whether, for Hermitian $H,K$, the value of
$\av(e^H,e^K)$ lies between the quantum partition function
$\tr e^{nH+mK}$ and $\tr(e^{nH}e^{mK})$ [1]. Cha and Lee disproved the
upper inequality in general [4]. Here the spectral multiplicities and
both word counts are arbitrary.

\paragraph{Theorem.}
If $A,B\in M_d(\mathbb C)$ are positive semidefinite and $A$ has at most
two distinct eigenvalues, then
\begin{equation}
 \tr[(A^{n/m}B)^m]\ \le\ \av(A,B)\ \le\ \tr(A^nB^m)
 \qquad(n,m\ge1).                                                    \label{eq:main}
\end{equation}
For $A=e^H$ and $B=e^K$, the first term dominates $\tr e^{nH+mK}$ by
the Araki--Lieb--Thirring inequality and the Lie--Trotter formula
[2, Lemma~2]. Interchanging the two factors gives the partition-function
and upper bounds also when $K$ has at most two spectral values.

\paragraph{Proof.}
Write $A=aP+bQ$ with $P,Q$ complementary orthogonal projections and
$a,b\ge0$. A scalar $A$ gives equality. Let
$\mathcal R=\{r\in\mathbb N_0^m\,:\,\sum_i r_i=n\}$, of size
$N=\binom{n+m-1}{m-1}$. By cyclicity [2, Lemma~3],
\[
 \av(A,B)=N^{-1}\sum_{r\in\mathcal R}
              \tr(A^{r_1}B A^{r_2}B\cdots A^{r_m}B).
\]
For example, marking one of the $m$ occurrences of $B$ in each word
counts each word $m$ times. Rotating the marked letter to the end leaves
a weak composition and a rotation offset in $\{0,\ldots,n+m-1\}$. The
trace is independent of that offset.

\newpage
For a sequence $E=(E_1,\ldots,E_m)\in\{P,Q\}^m$, set
$w(E)=\tr(E_1 B E_2 B\cdots E_m B)$. Let
$\lambda_i=a$ when $E_i=P$ and $\lambda_i=b$ when $E_i=Q$.
The scalar multiplying $w(E)$ in the average is
\[
 \kappa_k=N^{-1}\sum_{r\in\mathcal R}\prod_i\lambda_i^{r_i},
 \qquad k=\#\{i:E_i=Q\}.
\]
It depends only on $k$ because $\mathcal R$ is invariant under coordinate
permutations. With $c_k=\sum_{E:\#Q=k}w(E)$ we therefore have
$\av(A,B)=\sum_k c_k\kappa_k$. Terms with exponent zero equal one,
including at a zero eigenvalue.

\paragraph{Positivity of the grouped weights.}
Put $U=B^{1/2}PB^{1/2}$ and $V=B^{1/2}QB^{1/2}$. Both are positive
semidefinite, and cyclicity gives
\[
 p(t)=\tr[(U+tV)^m]=\tr[((P+tQ)B)^m]=\sum_{k=0}^m c_k t^k.
\]
The polynomial form of the ordinary BMV theorem gives $c_k\ge0$ [3].
This use of the theorem permits arbitrary ranks of $P$ and $Q$ and arbitrary
complex off-diagonal blocks of $B$.

\paragraph{Scalar bounds.}
Every coordinate $r_i$ has average $n/m$ over $\mathcal R$. The
arithmetic--geometric mean inequality over compositions gives the lower
bound below. For the upper bound, apply the weighted arithmetic--geometric
mean inequality to $\prod_i(\lambda_i^n)^{r_i/n}$ for each composition
and average:
\[
 (a^{m-k}b^k)^{n/m}\ \le\ \kappa_k\
 \le\ \frac{(m-k)a^n+kb^n}{m}.
\]
Continuity extends the lower bound from positive eigenvalues to zero.
After multiplication by $c_k$ and summation, expansion of
$(A^{n/m}B)^m$ gives the left side of \eqref{eq:main}. For the upper
side, $p(1)=\tr B^m$ and
$p'(1)=m\tr[V(U+V)^{m-1}]=m\tr(QB^m)$. Therefore
\[
 \sum_k c_k\frac{(m-k)a^n+kb^n}{m}
 =a^n\tr(PB^m)+b^n\tr(QB^m)=\tr(A^nB^m).\quad\square
\]

\paragraph{An elementary check through five letters $B$.}
In the splitting $P\oplus Q$, let $B$ have diagonal blocks $X,Y$ and
upper-right block $Z$. With $m\le5$, a cyclic sequence of labels switches
blocks zero, two or four times. With zero switches its weight is $\tr X^m$
or $\tr Y^m$. Two switches give $\tr(X^i ZY^j Z^*)$ for some
$i,j\ge0$. With four switches there is at most one diagonal step, giving
$\tr[(ZZ^*)^2]$, $\tr[X(ZZ^*)^2]$, or $\tr[Y(Z^*Z)^2]$, up to cyclic
rotation. These traces are nonnegative because $X,Y$ are positive
semidefinite and the trace of a product of two positive semidefinite
matrices is nonnegative. Thus in these low degrees one obtains the required
weight signs directly. If one eigenspace is one-dimensional, one can diagonalize the other
diagonal block and adjust basis phases to make all entries nonnegative.
The general theorem allows both eigenspaces to have arbitrary rank.

\newpage
\section*{Scope, sharpness and an independent finite reproduction}
The theorem is a direct corollary of the composition identity and ordinary
BMV coefficient positivity. Priority for this particular deduction has
not been established. Since $\mathcal A_{n,m}(A,B)=\mathcal A_{m,n}(B,A)$,
a violation of the upper inequality requires at least three spectral values
in \emph{each} factor. Cha and Lee's example [4] attains that threshold:
take $A=1\oplus xC$, $B=xC\oplus1$,
$C=\left(\begin{smallmatrix}1&-1\\-1&1\end{smallmatrix}\right)$ and
$x=10^{-3}$. Both spectra equal $\{0,2x,1\}$. At $n=m=5$ their example
has $\mathcal A_{5,5}(A,B)-\tr(A^5B^5)>0$ [4, Proposition~6].
Adding $\varepsilon I$ to both matrices preserves the strict polynomial
gap for sufficiently small $\varepsilon>0$. Both then have three distinct
positive eigenvalues and Hermitian logarithms. Thus the threshold also
holds for the original exponential inputs.
The lower bound for general spectra remains open.

The lower comparison at $m=4$ is covered by an existing
\href{https://github.com/the-omega-institute/trureturing/blob/c1938e29dc5ffac04e065367b179487fb4cd9744/Blueprint/D5/S3/Quantum/Analysis/RefinedBMV/SingularFourth/ZeroSpectrumProduct.md}{three-value containment theorem}.
At arbitrary word counts, the lower comparison for $A=tP$ follows from the
\href{https://arxiv.org/abs/2610.31851}{spectral-pinching inequality}, and a one-dimensional eigenspace permits the nonnegative-entry basis of [2, Proposition~1]. The
\href{https://github.com/shunyaekam/oo-research/blob/research/discovery/sources/bmv-block-priority.md}{priority audit} found no previous general deduction for two positive levels of unrestricted multiplicities and all $m\ge5$.

\paragraph{Edition note.} Before applying either scalar inequality, group paths by how often they visit the second eigenspace. The ordinary BMV theorem gives the sign of each group total even when individual path traces have no known sign. The appendix uses three-eigenvalue computations to reproduce cases of an existing lower bound.

\paragraph{Appendix.}
A public Lean proof [5] covers all positive integers $n$ with $m=3$
when $A$ has at most three distinct eigenvalues, in any finite dimension.
The following direct computation reproduces the subfamily $d=3$,
$11\le n\le36$ using H\"agele's sufficient criterion [2, Lemmas~5--6].
For $a\le b\le c$ the eigenvalues of $A$, put
\[
 F(u,v,w)=\sum_{i+j+k=n}u^iv^jw^k-\tbinom{n+2}{2}(uvw)^{n/3}.
\]
All exponents in the sum are nonnegative integers. Define
\begin{align*}
 X&=F(a,a,c),\qquad Y=F(b,b,c),\qquad Z=F(a,b,c),\\
 Q_0&=F(b,b,a)F(c,c,a)+F(a,a,b)F(c,c,b),\\
 D&=Z^2-XY,\qquad T=Z(X+Y)+2XY,\qquad L=Q_0T^2-D^2(X+Y)^2.
\end{align*}
Nonnegativity of $L$ suffices for the lower inequality in \eqref{eq:main}
for $m=3$ under these appendix hypotheses. Indeed $L$ is H\"agele's
polynomial multiplied by $\binom{n+2}{2}^6$ at $c=1$ and is homogeneous
of degree $6n$. Choose $s=1$ if $3\mid n$ and $s=3$ otherwise. For $a<b$,
divide the $s$th roots of all three eigenvalues by $b^{1/s}-a^{1/s}$.
After this positive rescaling their roots have the form $(x,x+1,x+1+t)$
with $x,t\ge0$. Since $sn/3$ is integral, the substituted expression for
$L$ is an integer polynomial in $x,t$. Exact multiplication over all
26 cases gives a total of 1,670,650 nonzero coefficients, each positive.
Homogeneity and continuity cover all eigenvalues, including coincidences.
At $n=12$, the average contains 455 words and the polynomial has 2,324
nonzero coefficients, the smallest 2,912.
\href{https://github.com/shunyaekam/oo-research/blob/research/checks/bmv_degrees.py}{Replay for degrees 11 through 36}
and \href{https://github.com/shunyaekam/oo-research/blob/research/checks/next_intake_bmv.py}{standard-library degree-twelve replay}.

\small
\begin{enumerate}\setlength{\itemsep}{1pt}
\item IQOQI Vienna, \href{https://oqp.iqoqi.oeaw.ac.at/refinement-of-the-bessis-moussa-villani-conjecture/}{Open Quantum Problems, problem 40}, posed by D. H\"agele.
\item D. H\"agele, \href{https://arxiv.org/html/2611.00129v1}{On the lower bound of the refined Bessis--Moussa--Villani conjecture}, arXiv:2611.00129v1 (2026), Sections 2--5.
\item E. H. Lieb and R. Seiringer, \href{https://arxiv.org/abs/math-ph/0210027}{Equivalent forms of the Bessis--Moussa--Villani conjecture} (2004). H. R. Stahl, \href{https://arxiv.org/abs/1107.4875}{Proof of the BMV conjecture} (2013).
\item H. Cha and J. Lee, \href{https://arxiv.org/abs/2603.19927v6}{One-parameter counterexamples to the refined Bessis--Moussa--Villani conjecture} (2026), version 6.
\item The Omega Institute, \href{https://github.com/the-omega-institute/trureturing/blob/b5ab5d2d61023ceb5440dbfe3e143c39337cd04d/D5/S3/Quantum/Analysis/RefinedBMVThreeSpectrumCertificate.lean}{Three-spectrum product and exponential lower bounds}, public Lean proof archive (2026).
\end{enumerate}
\end{document}
